\begin{ouverture}[algorithmique]{18}{Algorithmique et programmation par blocs}{Partie 5 · Programmer et réinvestir}
{Séquences · Variables · Conditions · Boucles · Événements}
\objectif{Lire, compléter et expliquer un script.}
\objectif{Utiliser variables, conditions et répétitions.}
\objectif{Programmer des constructions et des réactions à un événement.}
\prerequis{Calcul littéral, angles et repérage.}
\end{ouverture}
\pagestyle{otMain}\essentiel\fontsize{11}{15}\selectfont
\phantomsection\label{cours-18-1}
\essSection{1}{Décrire une suite d’actions}
\begin{definition}Un algorithme est une suite précise d’instructions permettant une tâche. Un programme le traduit dans un langage. Une variable est une mémoire nommée : « mettre x à 3 » remplace son contenu par 3 ; « ajouter 2 à x » augmente son contenu actuel de 2.\end{definition}
\begin{savoirfaire}[Lire un script]Repérer l’événement de départ, puis suivre les blocs dans l’ordre. Construire un tableau des valeurs des variables. Dans une boucle « répéter n fois », toutes les instructions à l’intérieur sont reprises à chaque tour. Les indentations et lignes « fin » du livre indiquent leur appartenance.\end{savoirfaire}
\label{fig-cours-18-1-script-01}\figureGeo{script-01}

\begin{attention}Une instruction qui modifie une variable utilise sa valeur au moment où elle est exécutée. Permuter deux blocs peut changer le résultat. On initialise les variables avant de les faire évoluer, pour ne pas dépendre d’une ancienne exécution.\end{attention}
\clearpage
\phantomsection\label{cours-18-2}
\essSection{2}{Choisir, répéter et réagir}
\begin{definition}Une condition est vraie ou fausse. Le bloc « si … alors … sinon » choisit la branche correspondante. Un événement, comme le clic sur le drapeau, une touche pressée ou un message reçu, déclenche un script. Une condition de contact permet à un lutin de réagir à un autre.\end{definition}
\begin{exemple}Pour tracer un carré, stylo baissé, répéter 4 fois « avancer de c ; tourner de $90^\circ$ ». Pour un triangle équilatéral, le lutin tourne de $120^\circ$, angle extérieur, et non de $60^\circ$. On indique ici le sens de rotation.\end{exemple}
\begin{savoirfaire}[Mettre au point]Tester des entrées habituelles et des cas limites, comme une valeur juste égale au seuil d’une condition. Vérifier le nombre de tours et la place de l’affichage. Pour un dessin, préciser le point initial, l’orientation et l’état du stylo. Le résultat observé doit être expliqué par les instructions.\end{savoirfaire}
\begin{attention}Les niveaux 1 et 2 des attendus sont travaillés ici : séquences, conditions, répétitions, événements, variables et interactions simples. Le défi final propose une recherche supplémentaire ; les scripts parallèles complexes et les boucles « jusqu’à » ne sont pas exigés pour ce parcours.\end{attention}

\clearpage\rubrique{Découverte}
\begin{blocLivre}{ex-18-1}\begin{exercice}[1]{Suivre les valeurs}
Exécuter le script. Quelle valeur affiche-t-il ?
\label{fig-ex-18-1}\figureGeo{script-01}
\end{exercice}\end{blocLivre}
\begin{blocLivre}{ex-18-2}\begin{exercice}[1]{Une entrée utilisateur}
Le script demande un nombre x. Donner le résultat pour 5, puis exprimer le résultat pour une entrée quelconque.
\label{fig-ex-18-2}\figureGeo{script-02}
\end{exercice}\end{blocLivre}
\begin{blocLivre}{ex-18-3}\begin{exercice}[1]{Une boucle d’ajouts}
Exécuter la boucle du script et donner l’affichage.
\label{fig-ex-18-3}\figureGeo{script-03}
\end{exercice}\end{blocLivre}
\begin{blocLivre}{ex-18-4}\begin{exercice}[1]{Une condition au seuil}
Le script classe un nombre avec le test x $\ge10$. Donner les messages pour x = 9 et x = 10.
\label{fig-ex-18-4}\figureGeo{script-04}
\end{exercice}\end{blocLivre}
\begin{blocLivre}{ex-18-5}\begin{exercice}[1]{Un produit répété}
Le script part de x = 3 et double x quatre fois. Donner les valeurs successives et l’affichage final.
\label{fig-ex-18-5}\figureGeo{script-05}
\end{exercice}\end{blocLivre}

\rubrique{Entraînement}
\begin{blocLivre}{ex-18-6}\begin{exercice}[2]{Une somme construite}
Le script utilise s et n. Donner les valeurs de s après chaque tour et expliquer le résultat.
\label{fig-ex-18-6}\figureGeo{script-06}
\end{exercice}\end{blocLivre}
\begin{blocLivre}{ex-18-7}\begin{exercice}[2]{Tracer un carré}
Le lutin part de (0;0), orienté vers la droite, stylo baissé. Décrire la figure et la position finale du script.
\label{fig-ex-18-7}\figureGeo{script-07}
\end{exercice}\end{blocLivre}
\begin{blocLivre}{ex-18-8}\begin{exercice}[2]{Tracer un rectangle}
Décrire le dessin du script, avec le même départ et le stylo baissé. Donner son périmètre.
\label{fig-ex-18-8}\figureGeo{script-08}
\end{exercice}\end{blocLivre}
\begin{blocLivre}{ex-18-9}\begin{exercice}[2]{Deux tests de parité}
Le bloc « reste de x divisé par 2 » est nul si x est pair. Que dit le script pour 18 et pour 19 ?
\label{fig-ex-18-9}\figureGeo{script-09}
\end{exercice}\end{blocLivre}
\begin{blocLivre}{ex-18-10}\begin{exercice}[2]{Utiliser une touche}
Le script est déclenché par la touche espace et ajoute 1 à score. score vaut 0 avant les appuis et n’est pas réinitialisé par ce script. Quelle valeur obtient-on après cinq appuis successifs ?
\label{fig-ex-18-10}\figureGeo{script-10}
\end{exercice}\end{blocLivre}

\rubrique{Problèmes}
\begin{blocLivre}{ex-18-11}\begin{exercice}[3]{Un triangle équilatéral}
Le lutin avance de 40 pixels et tourne de $120^\circ$, trois fois. Il part orienté à droite, stylo baissé. Justifier la figure et expliquer pourquoi $60^\circ$ ne convient pas.
\label{fig-ex-18-11}\figureGeo{script-11}
\end{exercice}\end{blocLivre}
\begin{blocLivre}{ex-18-12}\begin{exercice}[3]{Une interaction simple}
Le script teste si le lutin touche le lutin Guide. Quel message affiche-t-il selon que le contact est présent ou absent ? Pourquoi faut-il relancer ce script pour tester un nouveau contact ?
\label{fig-ex-18-12}\figureGeo{script-12}
\end{exercice}\end{blocLivre}
\begin{blocLivre}{ex-18-13}\begin{exercice}[3]{Une dimension saisie}
Le script demande une longueur positive c puis trace une figure. Donner le résultat pour c = 25 et exprimer la longueur totale tracée en fonction de c.
\label{fig-ex-18-13}\figureGeo{script-13}
\end{exercice}\end{blocLivre}
\begin{blocLivre}{ex-18-14}\begin{exercice}[3]{Chercher une erreur de boucle}
Un script met s à 0 à l’intérieur de la boucle, puis ajoute 3 à s, quatre fois. Il affiche s après la boucle. Expliquer le résultat et corriger pour obtenir 12.
\label{fig-ex-18-14}\figureGeo{script-14}
\end{exercice}\end{blocLivre}
\begin{blocLivre}{ex-18-15}\begin{exercice}[3]{Concevoir et vérifier}
Créer un script qui demande x, calcule $3x+5$, puis dit « positif » si le résultat est strictement positif et « nul ou négatif » sinon. Tester les entrées 2, $-2$ et $-5/3$.
\label{fig-ex-18-15}\figureGeo{script-15}
\end{exercice}\end{blocLivre}

\clearpage
